\SourceLeaf{245}%
\Needspace{6\baselineskip}\begin{center}\emph{[\textup{\textsc{Gerling}} an \textup{\textsc{Gauss}}. Marburg, 26. Februar 1844.]}\end{center}

\{\mbox{.\ .\ .\ .\ .\ .\ .} Die Notiz über das Buch von \textsc{Bolyai} habe ich bereits benutzt, um das Buch zu verschreiben. Unsere Bibliothek hat sich, mit diesem ausführlichen Titel ausgerüstet, unmittelbar an die Buchhändler wenden können. Gleichzeitig ist auch das Buch von \textsc{Lobatschewsky} aus Berlin verschrieben und so werden wir denn wohl hoffentlich nach nicht allzu langer Zeit in den Besitz kommen. Ich wiederhole meinen herzlichen Dank für Ihre gütige Bemühung. Die russischen Abhandlungen sind mir freilich unzugänglich, da ich selbst das Buchstabiren, was ich als Student gelernt, wieder vergessen habe. Den Aufsatz in \textsc{Crelle} Band 17 habe ich inzwischen auch nachgesehen und hängt darin allerdings alles an dem j'ai démontré. --- Das russische Steppenland scheint demnach doch ein geeigneter Boden für diese Speculationen, denn \textsc{Schweikart} (jetzt in Königsberg) ersann seine »Astral-Geometrie«, während er in Charkow war. \mbox{.\ .\ .\ .\ .\ .\ .}\}

\par\bigskip\Needspace{6\baselineskip}

\begin{center}\textsc{Gauss} an \textsc{Schumacher}. Göttingen, 28. November 1846.\end{center}

\mbox{.\ .\ .\ .\ .\ .\ .\ .\ .} Ich habe kürzlich Veranlassung gehabt, das Werkchen von \textsc{Lobatschewsky} (Geometrische Untersuchungen zur Theorie der Parallellinien. Berlin 1840, bei G. Fincke. 4 Bogen stark) wieder durchzusehen. Es enthält die Grundzüge derjenigen Geometrie, die Statt finden müsste und strenge consequent Statt finden könnte, wenn die Euklidische nicht die wahre ist. Ein gewisser \textsc{Schweikart}\HistoricalNote{*)}{*)}{Früher in Marburg, jetzt Professor der Jurisprudenz in Königsberg.} nannte eine solche Geometrie Astralgeometrie, \textsc{Lobatschewsky} imaginäre Geometrie. Sie wissen, dass ich schon seit 54 Jahren (seit 1792) dieselbe Überzeugung habe (mit einer gewissen spätern Erweiterung, deren ich hier nicht erwähnen will); materiell für mich Neues habe ich
